% Part: many-valued-logic
% Chapter: syntax-and-semantics
% Section: connectives

\documentclass[../../../include/open-logic-section]{subfiles}

\begin{document}

\olfileid{mvl}{syn}{con}

\olsection{भाषा आणि संयोजके}

अभिजात विधानीय तर्कशास्त्रात आणि इतर अनेक तर्कशास्त्रांत
\emph{विधानीय स्थिर} व \emph{संयोजक} यांचा संच वापरला जातो.
उदाहरणार्थ, पुढील चिन्हे आपण आदिम मानतो:
\begin{enumerate}
\tagitem{prvFalse}{!!{falsity} दर्शवणारा विधानीय स्थिर~$\lfalse$.}{}
\tagitem{prvTrue}{!!{truth} दर्शवणारा विधानीय स्थिर~$\ltrue$.}{}
\item तार्किक संयोजके:
  \startycommalist
  \iftag{prvNot}{\ycomma $\lnot$ (नकारण)}{}%
  \iftag{prvAnd}{\ycomma $\land$ (संधियोग)}{}%
  \iftag{prvOr}{\ycomma $\lor$ (विकल्पयोग)}{}%
  \iftag{prvIf}{\ycomma $\lif$ (!!{conditional})}{}%
  \iftag{prvIff}{\ycomma $\liff$ (!!{biconditional})}{}%
\end{enumerate}
\iftag{defNot,defOr,defAnd,defIf,defIff,defTrue,defFalse,defEx,defAll}{%
वरील आदिम संयोजकांव्यतिरिक्त संक्षिप्त रूपे म्हणून
परिभाषित केलेली चिन्हेही आपण वापरतो; उदाहरणार्थ:
\startycommalist
  \iftag{defNot}{\ycomma $\lnot$ (नकारण)}{}%
  \iftag{defAnd}{\ycomma $\land$ (संधियोग)}{}%
  \iftag{defOr}{\ycomma $\lor$ (विकल्पयोग)}{}%
  \iftag{defIf}{\ycomma $\lif$ (!!{conditional})}{}%
  \iftag{defIff}{\ycomma $\liff$ (!!{biconditional})}{}%
  \iftag{defFalse}{\ycomma $\lfalse$ (!!{falsity})}{}%
  \iftag{defTrue}{\ycomma $\ltrue$ (!!{truth})}.}{}

बहुमूल्य तर्कशास्त्रांतही हीच संयोजके वापरली जातात. तथापि,
उदाहरणार्थ संधियोगाच्या वेगवेगळ्या आवृत्त्या एकाच तर्कशास्त्रात
असणे अनेकदा उपयुक्त ठरते; त्या आवृत्त्या वेगळ्या दाखवण्यासाठी
भिन्न चिन्हे लागतात. काही बहुमूल्य तर्कशास्त्रांत अभिजात
तर्कशास्त्रात समतुल्य संयोजक नसलेली संयोजकेही असतात.
म्हणून आपण नेहमीपेक्षा थोडी अधिक सामान्य चौकट घेऊ.

\begin{defn}
  \emph{विधानीय भाषा} ही \emph{संयोजकां}च्या $\Lang L$
  या संचाने बनते. प्रत्येक संयोजक~$\star$ ची \emph{स्थानसंख्या}
  असते; स्थानसंख्या~$n$ असलेला संयोजक \emph{$n$-स्थानी}
  म्हणतात. स्थानसंख्या~$0$ असलेल्या संयोजकांना
  \emph{स्थिर} असेही म्हणतात; स्थानसंख्या~$1$ असलेल्या
  संयोजकांना \emph{एकस्थानी}, आणि स्थानसंख्या~$2$ असलेल्या
  संयोजकांना \emph{द्विस्थानी} म्हणतात.
\end{defn}

\begin{ex}
  विधानीय तर्कशास्त्राच्या मानक भाषेत~$\Lang L_0$ पुढील
  संयोजके (त्यांच्या स्थानसंख्यांसह) असतात:
  $\lfalse$~($0$),
  $\lnot$~($1$),
  $\land$~($2$),
  $\lor$~($2$),
  $\lif$~($2$). आपण विचारात घेणारी बहुतेक तर्कशास्त्रे
  ही भाषा वापरतात. काही तर्कशास्त्रे प्रस्थापित संकेतांनुसार
  काही संयोजकांसाठी वेगळी चिन्हे वापरतात. उदाहरणार्थ,
  गुणाकार तर्कशास्त्रात संधियोगाचे चिन्ह $\land$ ऐवजी
  अनेकदा~$\odot$ असते. कधी नवीन कारक जोडणेही सोयीचे
  ठरते; उदाहरणार्थ, $\triangle$ हा निश्चितता कारक
  ($1$-स्थानी).
\end{ex}

\end{document}
